it suffices, by the definition of @@M0@@ as left adjoint to the direct-image functor, to define an @@M1@@-linear homomorphism
@@M2@@
For this, we must give, for every open @@M3@@ of @@M4@@, a homomorphism:
@@M5@@
compatible with restrictions to the opens of @@M6@@, and take the corresponding homomorphism on the associated sheaves. Homomorphism (2.5.2) is defined as follows: let @@M7@@, where, @@M8@@ being of degree @@M9@@, @@M10@@; we then associate to @@M11@@ the section @@M12@@ of @@M13@@ over @@M14@@, writing @@M15@@ for the canonical homomorphism @@M16@@ (cf. EGA II 2.6.3 and 2.6.4). We easily check the necessary compatibilities, whence, by gluing, homomorphism (2.5.1).

It is clear that, in the case where @@M17@@ is quasi-coherent, the homomorphism thus constructed is the one in EGA II (3.3.4.1).

The homomorphism of functors @@M18@@ on the category of @@M19@@-Modules gives, by passage to the derived categories, a homomorphism of functors on @@M20@@:
\sgaStableTarget{sga6.e06.eq.2-5-2}@@M21@@

\sgaStableTarget{sga6.e06.prop.2-6}\textbf{Proposition 2.6.} \uline{Let} @@M22@@ \uline{be a quasi}-\uline{compact scheme}, @@M23@@ \uline{a locally free} @@M24@@\uline{-Module of finite type}, @@M25@@, \uline{and} @@M26@@ \uline{a perfect complex on} @@M27@@. \uline{Then} \uline{there exists an integer} @@M28@@ \uline{such that, for} @@M29@@, \uline{the following properties hold}:
